%% ===========================================================================
%%  「修士論文を書くにあたって」だけを 1 冊の PDF にするためのファイル。
%%
%%      latexmk guide-only.tex
%%
%%  main.tex とは独立した文書です。設定だけを同じ preamble.tex から読んでいるので、
%%  本編と同じ組版・同じ参考文献の形式になります。
%%  guide.tex を直せば、本編にもこちらにも反映されます。
%%
%%  GitHub Actions がこれを使って guide-only.pdf を作っています。
%% ===========================================================================

%% 配って読んでもらうものなので、白紙のページが入らないよう片面にする。
\documentclass[a4paper,11pt,oneside]{ltjsbook}

% !TEX root = main.tex

% ---------------------------------------------------------------------------
% このファイルは単体ではコンパイルできません（\documentclass がない、設定の断片です）。
% 間違えて latexmk preamble.tex としたときに、
% 何千行もエラーを出す前に止めるための仕掛けです。
% ---------------------------------------------------------------------------
\ifdefined\chapter\else
  \errmessage{^^J
   ================================================================^^J
   preamble.tex is a fragment. It cannot be compiled on its own.^^J
   preamble.tex は単体ではコンパイルできません。^^J
   latexmk                 (本編ぜんぶの PDF)^^J
   ================================================================^^J}
  \expandafter\stop
\fi
%% ===========================================================================
%%  パッケージと設定
%%
%%  ここは書き方の設定です。論文を書くのに触る必要はありません。
%%  自分用に書き換えるのは main.tex だけで済むようにしてあります。
%%
%%  余白を変えたい、参考文献の形式を変えたい、といったときだけここを開いてください。
%% ===========================================================================

%% ===========================================================================
%%  参考文献の方式（どちらか一方だけを有効にする）
%% ---------------------------------------------------------------------------
%%   biblatex + biber : .bib をそのまま使え、日本語文献も UTF-8 で素直に通る。
%%   natbib  + bst    : 天文分野で使われてきた aasjournalv7.bst をそのまま使う。
%%  本文の書き方（\citep{} / \citet{}）はどちらでも同じです。
%%  切り替えたら、いちど latexmk -C で中間ファイルを消してから作り直すこと。
%% ===========================================================================
\def\BibStyleBiblatex{}     % ← 既定
%\def\BibStyleNatbib{}      % ← こちらを使うときは上の行をコメントアウトする


%% ===========================================================================
%%  パッケージと設定
%% ===========================================================================

%% 共通プリアンブル（フォント・数式・図・単位・マクロ）。
%% 中身は hh-template.sty を見てください。全テンプレートで共通です。
%% 余白と hyperref はこの下で論文用に設定するので、ここでは外しておく。
\usepackage[nogeometry,nohyperref]{hh-template}

%% 版面。ltjsbook の既定は余白が広いので詰め、両面印刷でのど側を広めに取る。
\usepackage[a4paper,inner=30truemm,outer=22truemm,top=25truemm,bottom=25truemm]{geometry}

%% 図表
%% caption は ltjsbook を知らないので
%%   Package caption Warning: Unknown document class ...
%% という警告を必ず出しますが、実害はありません。無視してください。
\usepackage{caption}
\usepackage{subcaption}      % 図(a)(b) を並べる
\captionsetup{font=small,labelfont=bf}
\captionsetup[sub]{font=small,labelformat=parens,labelfont=bf,%
                   justification=raggedright,singlelinecheck=false}
\usepackage{longtable}       % 何ページにもわたる表
\usepackage{pdflscape}       % 横長の表を 90 度回して 1 ページに収める
\graphicspath{{fig/}}        % \includegraphics{foo} で fig/foo を探す

%% 目次
\usepackage[nottoc,notlot,notlof]{tocbibind}   % 引用文献を目次に載せる
\setcounter{tocdepth}{2}      % 目次に出すのは節まで
\setcounter{secnumdepth}{3}   % 番号を振るのは小節まで

%% 参考文献の実体
\ifdefined\BibStyleBiblatex
  \usepackage[backend=biber,
              style=authoryear,
              natbib=true,        % \citep \citet を使えるようにする
              maxcitenames=2,     % 3 人以上は「Hotta et al.」
              maxbibnames=99,     % 文献リストでは全員書く
              giveninits=true,    % 名は頭文字だけ
              uniquename=init,
              uniquelist=false,   % 3 人以上を必ず「Hotta et al.」に truncate する
              sorting=nyt,
              doi=false,isbn=false,url=false,eprint=false,
              backref=false]{biblatex}
  \addbibresource{thesis.bib}
\else
  \usepackage{natbib}
  \setcitestyle{authoryear,round,aysep={},citesep={;}}
\fi

%% リンク。LuaLaTeX なので日本語のしおりはそのまま通る（pxjahyper は不要）。
%% 紙で提出する版は下の行と入れ替えて文字を黒にすること。
\usepackage[unicode,colorlinks=true,allcolors=blue,pdfencoding=auto]{hyperref}
%\usepackage[unicode,colorlinks=false,pdfencoding=auto]{hyperref}  % 印刷用
\usepackage{url}

%% 表紙
\usepackage{thesiscover}

%% 添削のときだけ使うもの。行番号を振ると「12 行目のここ」と指摘しやすくなる。
%% 事務に提出する版では必ずコメントアウトすること。
%\usepackage{lineno}
%\linenumbers

%% この論文だけで使うマクロが必要ならここに書く。
%% ただし、数式の記法をマクロで置き換えるのはやめてください。
%% ここだけで通じる書き方を覚えても、他所で論文を書くときには使えません。


\begin{document}

\frontmatter
\thispagestyle{empty}
\vspace*{5\zw}
\begin{center}
  {\LARGE 修士論文を書くにあたって\par}
  \vspace{2\zw}
  {\large 名古屋大学宇宙地球環境研究所\par}
  \vspace{0.5\zw}
  {\large 堀田 英之\par}
  \vspace{1\zw}
  {\large \today\par}
\end{center}
\clearpage

\tableofcontents

\mainmatter
% !TEX root = main.tex
%% ===========================================================================
%%  執筆の手引き
%%
%%  論文の中身ではなく書き方の説明です。書き始める前に一度目を通してください。
%%  自分の原稿が形になってきたら、main.tex の % !TEX root = main.tex
%% ===========================================================================
%%  執筆の手引き
%%
%%  論文の中身ではなく書き方の説明です。書き始める前に一度目を通してください。
%%  自分の原稿が形になってきたら、main.tex の % !TEX root = main.tex
%% ===========================================================================
%%  執筆の手引き
%%
%%  論文の中身ではなく書き方の説明です。書き始める前に一度目を通してください。
%%  自分の原稿が形になってきたら、main.tex の \input{guide} を
%%  コメントアウトして構いません。
%% ===========================================================================

\chapter{執筆の手引き}
\label{chap:guide}

\section{文章の書き方}

\subsection{論文は主張の連なりである}

修士論文でよくある失敗は、やったことを時系列で並べてしまうことである。
読者が知りたいのは君が何を試したかの記録ではなく、
\textbf{何が新しく分かったのか}と\textbf{なぜそう言えるのか}の二つだけである。
まず結論を決め、その結論を支えるのに必要な材料だけを、必要な順に並べる。
やったが結論に効かなかった作業は、付録に回すか、思い切って落とす。

\subsection{一つの段落には一つの主張}

段落の先頭の一文（トピックセンテンス）にその段落の主張を書き、
残りの文でそれを支える。段落の頭だけを拾い読みして論旨が通るなら、
その章はよく書けている。

\subsection{日本語の作法}

\begin{description}
  \item[主語と述語を近づける] 間に修飾句を挟むほど読みにくくなる。
    一文が 3 行を超えたら、ほぼ確実に切れる。
  \item[「〜と考えられる」を減らす] 自分の結果について書くときは
    「〜である」と言い切る。断定できないなら、なぜ断定できないかを書く。
  \item[主観語を使わない] 「非常に」「かなり」「劇的に」は情報を持たない。
    「3 倍大きい」と書く。
  \item[「思う」は使わない] 論文は感想文ではない。
  \item[用語を統一する] 「対流セル」と「対流構造」を無意識に混ぜない。
    同じものは同じ語で呼び、違うものは違う語で呼ぶ。
  \item[数式も文の一部] 式の前後の助詞をつなげて音読し、
    日本語として成立するか確かめる。
\end{description}

\subsection{単位と有効数字}

物理量には必ず単位をつける。数値の有効数字は測定・計算の精度に見合った桁数にする。
計算機が出した 8 桁をそのまま書き写さない。
単位つきの量は \texttt{siunitx} を使うと組版が揃う。

\begin{verbatim}
太陽半径は \qty{6.96e8}{m}、自転角速度は \qty{2.6e-6}{\per\second} である。
速度は \qty{2}{\km\per\second}、範囲は \qtyrange{1}{5}{\km\per\second}。
\end{verbatim}

\noindent
従来どおり \verb|$2~\mathrm{km~s^{-1}}$| と手で書いてもよいが、
論文全体でどちらかに統一すること。

\subsection{書き始められないとき}

完成形を頭の中で作ってから書こうとすると、いつまでも書き始められない。
図を先に作り、それを並べ替えて話の順番を決め、
各図に対して説明の段落を書く、という順で進めるとよい。
文章の質は推敲で上げるもので、初稿で上げるものではない。


\section{\LaTeX{} の使い方}
\label{sec:guide_latex}

\subsection{このテンプレートの構成}

\begin{description}
  \item[\texttt{main.tex}] プリアンブル、表紙に載る情報、概要、章の並び。
    自分用に書き換えるのは主にこのファイルの上半分である。
  \item[\texttt{body.tex}] 本文。序論から結論まで。
  \item[\texttt{guide.tex}] この手引き。書き終わったら外してよい。
  \item[\texttt{thesis.bib}] 引用文献のデータベース。
  \item[\texttt{achilles.sty}] 全テンプレート共通のプリアンブル。
    原本はリポジトリの \texttt{common/} にある。
  \item[\texttt{thesiscover.sty}] 表紙のレイアウト。通常は触らない。
  \item[\texttt{fig/}] 図を置くところ。
\end{description}

\subsection{コンパイル}

このディレクトリで
\begin{verbatim}
latexmk
\end{verbatim}
とすれば \texttt{main.pdf} ができる。書きながら自動で作り直したいときは
\begin{verbatim}
latexmk -pvc
\end{verbatim}
とする。中間ファイルがおかしくなったら \verb|latexmk -C| で全部消してやり直す。
エンジンは LuaLaTeX である。\texttt{platex} や \texttt{uplatex} では通らない。

\subsection{相互参照}

図・表・式・章には \verb|\label{}| をつけ、本文からは \verb|\ref{}| で参照する。
番号を手で書いてはいけない。あとで章を入れ替えたときに必ず破綻する。
ラベルには \verb|fig:|、\verb|tab:|、\verb|eq:|、\verb|sec:|、\verb|chap:| の
接頭辞をつけると、何を指しているか分かりやすい。

\begin{verbatim}
図\ref{fig:sample}に示すように……
式\eqref{eq:induction}より……
\end{verbatim}

\noindent
式の参照は \verb|\ref| ではなく \verb|\eqref| を使うと括弧がつく。

\subsection{引用}

\texttt{thesis.bib} に文献情報を書き、本文から \verb|\citep| か \verb|\citet| で引く。
NASA ADS の \textsf{Export Citation} から \textsf{BibTeX} を選べば、
そのまま貼り付けられる形式で得られる。手で打ち直さないこと。

\begin{verbatim}
小スケールダイナモが有効粘性を下げる \citep{hotta2021}。
\citet{hotta2021} はこれを高解像度計算で示した。
\end{verbatim}

\noindent
このテンプレートは \texttt{biblatex + biber} と \texttt{natbib + bst} の
どちらでも動くようにしてある。\texttt{main.tex} の冒頭で切り替える。
本文の書き方は両者で共通なので、あとから変えても原稿は書き直さなくてよい。

\subsection{数式}

番号をつけて後で参照する式は \texttt{align}、参照しない式は \texttt{align*} を使う。
\verb|$$ ... $$| は使わない（\LaTeX{} では非推奨で、前後の空きがおかしくなる）。
ベクトルは \verb|\vect{B}| と書くと太字斜体になる。
\verb|\grad|、\verb|\divergence|、\verb|\curl|、\verb|\pdif{}{}| などの
短縮マクロを \texttt{achilles.sty} に用意してある。

\subsection{よくあるつまずき}

\begin{description}
  \item[エラーが出たら一番上から読む] 後ろのエラーは前のエラーの巻き添えである
    ことが多い。\verb|! LaTeX Error:| で始まる最初の行と、その直前の行番号を見る。
  \item[Undefined control sequence] マクロの綴りが違うか、
    必要なパッケージを読み込んでいない。
  \item[Overfull \textbackslash hbox] 行がはみ出している。長い英単語や URL が原因。
    数 pt なら無視してよいが、目で見て飛び出していたら直す。
  \item[番号が \texttt{??} になる] 参照が解決していない。もう一度コンパイルする
    （\texttt{latexmk} なら自動でやる）。それでも直らないならラベルの綴り間違い。
  \item[図が出ない] ファイル名か拡張子が違う。LuaLaTeX が読めるのは
    PDF・PNG・JPEG である。EPS は事前に PDF へ変換しておく。
\end{description}


\section{図表の作り方}
\label{sec:guide_figure}

審査員は本文より先に図を見る。図が分かりやすければ論文は半分伝わったも同然である。
逆に、図が読めない論文は本文がどれだけ良くても評価されない。

\subsection{形式}

\begin{description}
  \item[線画はベクタ形式で] 折れ線・等高線・散布図は PDF で保存する。
    拡大しても線が滑らかで、印刷しても綺麗である。
    Matplotlib なら \verb|plt.savefig('foo.pdf')| とするだけである。
  \item[画像はラスタ形式で] 太陽表面の強度分布のような画素の塊は PNG にする。
    PDF にすると数十 MB になり、PDF ビューアが固まる。
    \verb|dpi=300| 程度あれば印刷に耐える。
  \item[JPEG は使わない] 圧縮による滲みが出る。写真以外では使わない。
\end{description}

\subsection{文字の大きさ}

図の中の文字が本文より小さいと読めない。
論文に貼ったときに\textbf{本文と同じかやや小さい程度}になるよう、
図を作る段階でフォントサイズを決める。
\verb|\includegraphics[width=0.5\linewidth]| のように後から縮めると
文字も一緒に縮むので、その分を見越しておく。
出来上がった PDF を 100\% 表示にして、読めるか自分の目で確かめること。

\subsection{色}

\begin{description}
  \item[虹色（jet）を使わない] 明度が単調でないため、
    データにない偽の構造が見えてしまう。Matplotlib の既定である
    \texttt{viridis} や、正負のある量には \texttt{RdBu\_r} のような
    発散型カラーマップを使う。
  \item[色覚多様性に配慮する] 日本人男性の約 5\% は赤と緑の区別がつきにくい。
    赤と緑だけで系列を区別しない。線種（実線・破線・点線）や
    マーカーの形も併せて変える。
  \item[白黒で印刷されても読めるか] 製本された論文が白黒コピーされることはある。
    色を消しても情報が失われないようにしておくと安全である。
  \item[色に意味を持たせたら統一する] 論文全体で「観測は黒、計算は赤」と決めたら、
    最後まで変えない。
\end{description}

\subsection{キャプション}

キャプションだけを読んで図が理解できるように書く。含めるべきものは
\begin{enumerate}
  \item 何を示した図か
  \item 軸の量と単位
  \item 線・色・記号がそれぞれ何を表すか
  \item どの計算・観測のデータか
\end{enumerate}
である。「結果」だけのキャプションは論外である。
図のキャプションは図の\textbf{下}、表のキャプションは表の\textbf{上}に置く。

\subsection{図を並べる}

\texttt{subcaption} パッケージで (a) (b) と番号を振れる。

\begin{verbatim}
\begin{figure}[tb]
  \centering
  \begin{subfigure}{0.48\linewidth}
    \includegraphics[width=\linewidth]{left.pdf}
    \caption{低解像度}\label{fig:comp_low}
  \end{subfigure}\hfill
  \begin{subfigure}{0.48\linewidth}
    \includegraphics[width=\linewidth]{right.pdf}
    \caption{高解像度}\label{fig:comp_high}
  \end{subfigure}
  \caption{解像度による対流構造の違い。}
  \label{fig:comp}
\end{figure}
\end{verbatim}

\subsection{図の位置が思いどおりにならないとき}

\LaTeX{} は図を「浮きもの」として扱い、ページの切れ目を見て適当な場所に置く。
\verb|[tb]| は「ページの上か下に置いてよい」という指定である。
\verb|[h]|（ここに置け）を多用すると、かえって図が後ろに溜まる。
最終稿になるまでは位置を気にせず、内容を書くことに集中したほうがよい。


\section{引用と剽窃}
\label{sec:guide_ethics}

これは作法の問題ではなく、研究者としての資格の問題である。
剽窃が発覚すれば学位は取り消される。指導教員にも累が及ぶ。
知らなかったでは済まないので、必ず読むこと。

\subsection{何が剽窃か}

他人の文章・図・アイデアを、出典を示さずに自分のものとして提示することは
すべて剽窃である。以下はいずれも剽窃にあたる。

\begin{itemize}
  \item 論文や web の文章をコピーして貼り、少し語尾を変えて自分の文にする
  \item 他人の図をキャプションに出典を書かずに載せる
  \item 先輩の修士論文から章をまるごと持ってくる
  \item 生成 AI が出力した文章を、出典も検証もなくそのまま載せる
  \item 自分の過去の論文から、断りなく文章を再利用する（自己剽窃）
\end{itemize}

\subsection{正しいやり方}

\begin{description}
  \item[事実や知見を借りるとき] 自分の言葉で書き直し、引用をつける。
    「対流層の底は $0.71\,R_\odot$ にある \citep{christensen-dalsgaard1991}。」
  \item[文章をそのまま使うとき] 引用符でくくり、出典を示す。
    ただし自然科学の論文で原文引用が必要になることはまずない。
    必要だと思ったら、たいていは自分の言葉で書き直すべき場面である。
  \item[図を借りるとき] キャプションの末尾に出典を書く。
    「\citet{hotta2021} の Figure 3 を改変。」
    商業出版社の論文図は、多くの場合、著者本人でも転載に許諾が要る。
    論文の web ページにある \textsf{Permissions} や \textsf{RightsLink} から
    手続きする。許諾が取れないなら、自分で描き直す。
\end{description}

\noindent
どこまで自分の言葉にすればよいかの基準は、
語順を入れ替えたか同義語に置き換えたかではない。
\textbf{元の文章を見ずに、理解した内容を自分で書けるか}である。
見ないと書けないなら、まだ理解していない。

\subsection{生成 AI の扱い}

英文の校正や、書き出しに詰まったときの相談に使うのは構わない。
ただし次の 3 点は守ること。

\begin{enumerate}
  \item 出力された内容の正しさは自分で確認する。
    存在しない論文を引用してくることがある。
  \item 引用文献は必ず ADS などの一次情報で実在を確かめる。
  \item 使い方について、指導教員に隠さない。
    大学や学会が定める方針にも従うこと。
\end{enumerate}

\noindent
自分が説明できない文章を論文に載せてはいけない。
審査で問われるのは論文ではなく、書いた本人だからである。

\subsection{データと図の加工}

都合の悪い点を消す、エラーバーを小さく見せる、
軸を切って差を大きく見せる、といった操作は捏造・改竄である。
処理を加えたなら、加えたことを本文に書く。
「見やすさのために平滑化した」と書けば済む話を、黙ってやると不正になる。


\section{原稿の管理と添削の受け方}
\label{sec:guide_review}

\subsection{Git で履歴を残す}

\texttt{.tex} はテキストファイルなので Git で管理できる。
締切前に「昨日の版に戻したい」と言い出す前に、最初から履歴を残しておくこと。
節を書き終えるたび、あるいは一日の終わりに
\begin{verbatim}
git add -A
git commit -m "序論の先行研究を追加"
\end{verbatim}
としておけばよい。中間ファイルや PDF は \texttt{.gitignore} で除外してある。
バックアップは Git だけに頼らず、クラウド（大学のストレージなど）にも置くこと。

\subsection{添削に出す前に自分でやること}

指導教員の時間を、自分で見つけられる誤りに使わせないこと。
出す前に次を済ませておく。

\begin{itemize}
  \item \verb|\todo| と \verb|\red| が残っていないか検索する
  \item \texttt{??} や \texttt{[?]} が PDF に出ていないか確認する
  \item 図がすべて本文から参照されているか確認する
  \item 通しで一度、声に出して読む（黙読よりずっと効く。文の切れ目が分かる）
  \item 数式の記号がすべて本文で定義されているか確認する
\end{itemize}

\subsection{添削の受け方}

\begin{description}
  \item[早く出す] 完璧になってから見せようとすると間に合わない。
    構成だけの段階、章ごと、と何度も出したほうが結果的に速い。
  \item[版に名前をつける] PDF は \texttt{thesis\_20270105.pdf} のように
    日付をつけて渡す。「最新版」が何を指すか分からなくなるのを防ぐ。
  \item[指摘の理由を理解する] 直せと言われた箇所だけ直すと、
    同じ誤りを他の場所に残したまま次の版を出すことになる。
    なぜそう指摘されたのかを考え、同種の箇所をまとめて直す。
  \item[納得できないときは聞く] 指摘が間違っていることもある。
    黙って従うのも、黙って無視するのも良くない。
  \item[直したかどうかを返す] どの指摘をどう直したか、
    直さなかったならなぜかを添えて返すと、次の添削が速くなる。
\end{description}

\subsection{変更点を見せる}

\texttt{latexdiff} を使うと、前の版からの変更箇所に下線と取り消し線がついた
PDF を作れる。どこを直したかが一目で分かるので、
再添削のときはこれを添えると親切である。

\begin{verbatim}
latexdiff-vc --git -r HEAD~1 --flatten main.tex
\end{verbatim}

\subsection{提出前の最終確認}

\begin{itemize}
  \item 表紙の年度・題目・氏名・学籍番号・所属が学務の書類と一致しているか
  \item ページ番号が通っているか
  \item 目次・図目次・表目次の番号が本文と合っているか
  \item 引用文献が本文の引用とすべて対応しているか
  \item 手引き（この章）を \texttt{main.tex} から外したか
  \item 行番号（\texttt{lineno}）を切ったか
  \item 印刷して提出するなら、リンクの色を黒にしたか
  \item PDF を実際に印刷して、図の文字が読めるか確かめたか
\end{itemize}
 を
%%  コメントアウトして構いません。
%% ===========================================================================

\chapter{執筆の手引き}
\label{chap:guide}

\section{文章の書き方}

\subsection{論文は主張の連なりである}

修士論文でよくある失敗は、やったことを時系列で並べてしまうことである。
読者が知りたいのは君が何を試したかの記録ではなく、
\textbf{何が新しく分かったのか}と\textbf{なぜそう言えるのか}の二つだけである。
まず結論を決め、その結論を支えるのに必要な材料だけを、必要な順に並べる。
やったが結論に効かなかった作業は、付録に回すか、思い切って落とす。

\subsection{一つの段落には一つの主張}

段落の先頭の一文（トピックセンテンス）にその段落の主張を書き、
残りの文でそれを支える。段落の頭だけを拾い読みして論旨が通るなら、
その章はよく書けている。

\subsection{日本語の作法}

\begin{description}
  \item[主語と述語を近づける] 間に修飾句を挟むほど読みにくくなる。
    一文が 3 行を超えたら、ほぼ確実に切れる。
  \item[「〜と考えられる」を減らす] 自分の結果について書くときは
    「〜である」と言い切る。断定できないなら、なぜ断定できないかを書く。
  \item[主観語を使わない] 「非常に」「かなり」「劇的に」は情報を持たない。
    「3 倍大きい」と書く。
  \item[「思う」は使わない] 論文は感想文ではない。
  \item[用語を統一する] 「対流セル」と「対流構造」を無意識に混ぜない。
    同じものは同じ語で呼び、違うものは違う語で呼ぶ。
  \item[数式も文の一部] 式の前後の助詞をつなげて音読し、
    日本語として成立するか確かめる。
\end{description}

\subsection{単位と有効数字}

物理量には必ず単位をつける。数値の有効数字は測定・計算の精度に見合った桁数にする。
計算機が出した 8 桁をそのまま書き写さない。
単位つきの量は \texttt{siunitx} を使うと組版が揃う。

\begin{verbatim}
太陽半径は \qty{6.96e8}{m}、自転角速度は \qty{2.6e-6}{\per\second} である。
速度は \qty{2}{\km\per\second}、範囲は \qtyrange{1}{5}{\km\per\second}。
\end{verbatim}

\noindent
従来どおり \verb|$2~\mathrm{km~s^{-1}}$| と手で書いてもよいが、
論文全体でどちらかに統一すること。

\subsection{書き始められないとき}

完成形を頭の中で作ってから書こうとすると、いつまでも書き始められない。
図を先に作り、それを並べ替えて話の順番を決め、
各図に対して説明の段落を書く、という順で進めるとよい。
文章の質は推敲で上げるもので、初稿で上げるものではない。


\section{\LaTeX{} の使い方}
\label{sec:guide_latex}

\subsection{このテンプレートの構成}

\begin{description}
  \item[\texttt{main.tex}] プリアンブル、表紙に載る情報、概要、章の並び。
    自分用に書き換えるのは主にこのファイルの上半分である。
  \item[\texttt{body.tex}] 本文。序論から結論まで。
  \item[\texttt{guide.tex}] この手引き。書き終わったら外してよい。
  \item[\texttt{thesis.bib}] 引用文献のデータベース。
  \item[\texttt{achilles.sty}] 全テンプレート共通のプリアンブル。
    原本はリポジトリの \texttt{common/} にある。
  \item[\texttt{thesiscover.sty}] 表紙のレイアウト。通常は触らない。
  \item[\texttt{fig/}] 図を置くところ。
\end{description}

\subsection{コンパイル}

このディレクトリで
\begin{verbatim}
latexmk
\end{verbatim}
とすれば \texttt{main.pdf} ができる。書きながら自動で作り直したいときは
\begin{verbatim}
latexmk -pvc
\end{verbatim}
とする。中間ファイルがおかしくなったら \verb|latexmk -C| で全部消してやり直す。
エンジンは LuaLaTeX である。\texttt{platex} や \texttt{uplatex} では通らない。

\subsection{相互参照}

図・表・式・章には \verb|\label{}| をつけ、本文からは \verb|\ref{}| で参照する。
番号を手で書いてはいけない。あとで章を入れ替えたときに必ず破綻する。
ラベルには \verb|fig:|、\verb|tab:|、\verb|eq:|、\verb|sec:|、\verb|chap:| の
接頭辞をつけると、何を指しているか分かりやすい。

\begin{verbatim}
図\ref{fig:sample}に示すように……
式\eqref{eq:induction}より……
\end{verbatim}

\noindent
式の参照は \verb|\ref| ではなく \verb|\eqref| を使うと括弧がつく。

\subsection{引用}

\texttt{thesis.bib} に文献情報を書き、本文から \verb|\citep| か \verb|\citet| で引く。
NASA ADS の \textsf{Export Citation} から \textsf{BibTeX} を選べば、
そのまま貼り付けられる形式で得られる。手で打ち直さないこと。

\begin{verbatim}
小スケールダイナモが有効粘性を下げる \citep{hotta2021}。
\citet{hotta2021} はこれを高解像度計算で示した。
\end{verbatim}

\noindent
このテンプレートは \texttt{biblatex + biber} と \texttt{natbib + bst} の
どちらでも動くようにしてある。\texttt{main.tex} の冒頭で切り替える。
本文の書き方は両者で共通なので、あとから変えても原稿は書き直さなくてよい。

\subsection{数式}

番号をつけて後で参照する式は \texttt{align}、参照しない式は \texttt{align*} を使う。
\verb|$$ ... $$| は使わない（\LaTeX{} では非推奨で、前後の空きがおかしくなる）。
ベクトルは \verb|\vect{B}| と書くと太字斜体になる。
\verb|\grad|、\verb|\divergence|、\verb|\curl|、\verb|\pdif{}{}| などの
短縮マクロを \texttt{achilles.sty} に用意してある。

\subsection{よくあるつまずき}

\begin{description}
  \item[エラーが出たら一番上から読む] 後ろのエラーは前のエラーの巻き添えである
    ことが多い。\verb|! LaTeX Error:| で始まる最初の行と、その直前の行番号を見る。
  \item[Undefined control sequence] マクロの綴りが違うか、
    必要なパッケージを読み込んでいない。
  \item[Overfull \textbackslash hbox] 行がはみ出している。長い英単語や URL が原因。
    数 pt なら無視してよいが、目で見て飛び出していたら直す。
  \item[番号が \texttt{??} になる] 参照が解決していない。もう一度コンパイルする
    （\texttt{latexmk} なら自動でやる）。それでも直らないならラベルの綴り間違い。
  \item[図が出ない] ファイル名か拡張子が違う。LuaLaTeX が読めるのは
    PDF・PNG・JPEG である。EPS は事前に PDF へ変換しておく。
\end{description}


\section{図表の作り方}
\label{sec:guide_figure}

審査員は本文より先に図を見る。図が分かりやすければ論文は半分伝わったも同然である。
逆に、図が読めない論文は本文がどれだけ良くても評価されない。

\subsection{形式}

\begin{description}
  \item[線画はベクタ形式で] 折れ線・等高線・散布図は PDF で保存する。
    拡大しても線が滑らかで、印刷しても綺麗である。
    Matplotlib なら \verb|plt.savefig('foo.pdf')| とするだけである。
  \item[画像はラスタ形式で] 太陽表面の強度分布のような画素の塊は PNG にする。
    PDF にすると数十 MB になり、PDF ビューアが固まる。
    \verb|dpi=300| 程度あれば印刷に耐える。
  \item[JPEG は使わない] 圧縮による滲みが出る。写真以外では使わない。
\end{description}

\subsection{文字の大きさ}

図の中の文字が本文より小さいと読めない。
論文に貼ったときに\textbf{本文と同じかやや小さい程度}になるよう、
図を作る段階でフォントサイズを決める。
\verb|\includegraphics[width=0.5\linewidth]| のように後から縮めると
文字も一緒に縮むので、その分を見越しておく。
出来上がった PDF を 100\% 表示にして、読めるか自分の目で確かめること。

\subsection{色}

\begin{description}
  \item[虹色（jet）を使わない] 明度が単調でないため、
    データにない偽の構造が見えてしまう。Matplotlib の既定である
    \texttt{viridis} や、正負のある量には \texttt{RdBu\_r} のような
    発散型カラーマップを使う。
  \item[色覚多様性に配慮する] 日本人男性の約 5\% は赤と緑の区別がつきにくい。
    赤と緑だけで系列を区別しない。線種（実線・破線・点線）や
    マーカーの形も併せて変える。
  \item[白黒で印刷されても読めるか] 製本された論文が白黒コピーされることはある。
    色を消しても情報が失われないようにしておくと安全である。
  \item[色に意味を持たせたら統一する] 論文全体で「観測は黒、計算は赤」と決めたら、
    最後まで変えない。
\end{description}

\subsection{キャプション}

キャプションだけを読んで図が理解できるように書く。含めるべきものは
\begin{enumerate}
  \item 何を示した図か
  \item 軸の量と単位
  \item 線・色・記号がそれぞれ何を表すか
  \item どの計算・観測のデータか
\end{enumerate}
である。「結果」だけのキャプションは論外である。
図のキャプションは図の\textbf{下}、表のキャプションは表の\textbf{上}に置く。

\subsection{図を並べる}

\texttt{subcaption} パッケージで (a) (b) と番号を振れる。

\begin{verbatim}
\begin{figure}[tb]
  \centering
  \begin{subfigure}{0.48\linewidth}
    \includegraphics[width=\linewidth]{left.pdf}
    \caption{低解像度}\label{fig:comp_low}
  \end{subfigure}\hfill
  \begin{subfigure}{0.48\linewidth}
    \includegraphics[width=\linewidth]{right.pdf}
    \caption{高解像度}\label{fig:comp_high}
  \end{subfigure}
  \caption{解像度による対流構造の違い。}
  \label{fig:comp}
\end{figure}
\end{verbatim}

\subsection{図の位置が思いどおりにならないとき}

\LaTeX{} は図を「浮きもの」として扱い、ページの切れ目を見て適当な場所に置く。
\verb|[tb]| は「ページの上か下に置いてよい」という指定である。
\verb|[h]|（ここに置け）を多用すると、かえって図が後ろに溜まる。
最終稿になるまでは位置を気にせず、内容を書くことに集中したほうがよい。


\section{引用と剽窃}
\label{sec:guide_ethics}

これは作法の問題ではなく、研究者としての資格の問題である。
剽窃が発覚すれば学位は取り消される。指導教員にも累が及ぶ。
知らなかったでは済まないので、必ず読むこと。

\subsection{何が剽窃か}

他人の文章・図・アイデアを、出典を示さずに自分のものとして提示することは
すべて剽窃である。以下はいずれも剽窃にあたる。

\begin{itemize}
  \item 論文や web の文章をコピーして貼り、少し語尾を変えて自分の文にする
  \item 他人の図をキャプションに出典を書かずに載せる
  \item 先輩の修士論文から章をまるごと持ってくる
  \item 生成 AI が出力した文章を、出典も検証もなくそのまま載せる
  \item 自分の過去の論文から、断りなく文章を再利用する（自己剽窃）
\end{itemize}

\subsection{正しいやり方}

\begin{description}
  \item[事実や知見を借りるとき] 自分の言葉で書き直し、引用をつける。
    「対流層の底は $0.71\,R_\odot$ にある \citep{christensen-dalsgaard1991}。」
  \item[文章をそのまま使うとき] 引用符でくくり、出典を示す。
    ただし自然科学の論文で原文引用が必要になることはまずない。
    必要だと思ったら、たいていは自分の言葉で書き直すべき場面である。
  \item[図を借りるとき] キャプションの末尾に出典を書く。
    「\citet{hotta2021} の Figure 3 を改変。」
    商業出版社の論文図は、多くの場合、著者本人でも転載に許諾が要る。
    論文の web ページにある \textsf{Permissions} や \textsf{RightsLink} から
    手続きする。許諾が取れないなら、自分で描き直す。
\end{description}

\noindent
どこまで自分の言葉にすればよいかの基準は、
語順を入れ替えたか同義語に置き換えたかではない。
\textbf{元の文章を見ずに、理解した内容を自分で書けるか}である。
見ないと書けないなら、まだ理解していない。

\subsection{生成 AI の扱い}

英文の校正や、書き出しに詰まったときの相談に使うのは構わない。
ただし次の 3 点は守ること。

\begin{enumerate}
  \item 出力された内容の正しさは自分で確認する。
    存在しない論文を引用してくることがある。
  \item 引用文献は必ず ADS などの一次情報で実在を確かめる。
  \item 使い方について、指導教員に隠さない。
    大学や学会が定める方針にも従うこと。
\end{enumerate}

\noindent
自分が説明できない文章を論文に載せてはいけない。
審査で問われるのは論文ではなく、書いた本人だからである。

\subsection{データと図の加工}

都合の悪い点を消す、エラーバーを小さく見せる、
軸を切って差を大きく見せる、といった操作は捏造・改竄である。
処理を加えたなら、加えたことを本文に書く。
「見やすさのために平滑化した」と書けば済む話を、黙ってやると不正になる。


\section{原稿の管理と添削の受け方}
\label{sec:guide_review}

\subsection{Git で履歴を残す}

\texttt{.tex} はテキストファイルなので Git で管理できる。
締切前に「昨日の版に戻したい」と言い出す前に、最初から履歴を残しておくこと。
節を書き終えるたび、あるいは一日の終わりに
\begin{verbatim}
git add -A
git commit -m "序論の先行研究を追加"
\end{verbatim}
としておけばよい。中間ファイルや PDF は \texttt{.gitignore} で除外してある。
バックアップは Git だけに頼らず、クラウド（大学のストレージなど）にも置くこと。

\subsection{添削に出す前に自分でやること}

指導教員の時間を、自分で見つけられる誤りに使わせないこと。
出す前に次を済ませておく。

\begin{itemize}
  \item \verb|\todo| と \verb|\red| が残っていないか検索する
  \item \texttt{??} や \texttt{[?]} が PDF に出ていないか確認する
  \item 図がすべて本文から参照されているか確認する
  \item 通しで一度、声に出して読む（黙読よりずっと効く。文の切れ目が分かる）
  \item 数式の記号がすべて本文で定義されているか確認する
\end{itemize}

\subsection{添削の受け方}

\begin{description}
  \item[早く出す] 完璧になってから見せようとすると間に合わない。
    構成だけの段階、章ごと、と何度も出したほうが結果的に速い。
  \item[版に名前をつける] PDF は \texttt{thesis\_20270105.pdf} のように
    日付をつけて渡す。「最新版」が何を指すか分からなくなるのを防ぐ。
  \item[指摘の理由を理解する] 直せと言われた箇所だけ直すと、
    同じ誤りを他の場所に残したまま次の版を出すことになる。
    なぜそう指摘されたのかを考え、同種の箇所をまとめて直す。
  \item[納得できないときは聞く] 指摘が間違っていることもある。
    黙って従うのも、黙って無視するのも良くない。
  \item[直したかどうかを返す] どの指摘をどう直したか、
    直さなかったならなぜかを添えて返すと、次の添削が速くなる。
\end{description}

\subsection{変更点を見せる}

\texttt{latexdiff} を使うと、前の版からの変更箇所に下線と取り消し線がついた
PDF を作れる。どこを直したかが一目で分かるので、
再添削のときはこれを添えると親切である。

\begin{verbatim}
latexdiff-vc --git -r HEAD~1 --flatten main.tex
\end{verbatim}

\subsection{提出前の最終確認}

\begin{itemize}
  \item 表紙の年度・題目・氏名・学籍番号・所属が学務の書類と一致しているか
  \item ページ番号が通っているか
  \item 目次・図目次・表目次の番号が本文と合っているか
  \item 引用文献が本文の引用とすべて対応しているか
  \item 手引き（この章）を \texttt{main.tex} から外したか
  \item 行番号（\texttt{lineno}）を切ったか
  \item 印刷して提出するなら、リンクの色を黒にしたか
  \item PDF を実際に印刷して、図の文字が読めるか確かめたか
\end{itemize}
 を
%%  コメントアウトして構いません。
%% ===========================================================================

\chapter{執筆の手引き}
\label{chap:guide}

\section{文章の書き方}

\subsection{論文は主張の連なりである}

修士論文でよくある失敗は、やったことを時系列で並べてしまうことである。
読者が知りたいのは君が何を試したかの記録ではなく、
\textbf{何が新しく分かったのか}と\textbf{なぜそう言えるのか}の二つだけである。
まず結論を決め、その結論を支えるのに必要な材料だけを、必要な順に並べる。
やったが結論に効かなかった作業は、付録に回すか、思い切って落とす。

\subsection{一つの段落には一つの主張}

段落の先頭の一文（トピックセンテンス）にその段落の主張を書き、
残りの文でそれを支える。段落の頭だけを拾い読みして論旨が通るなら、
その章はよく書けている。

\subsection{日本語の作法}

\begin{description}
  \item[主語と述語を近づける] 間に修飾句を挟むほど読みにくくなる。
    一文が 3 行を超えたら、ほぼ確実に切れる。
  \item[「〜と考えられる」を減らす] 自分の結果について書くときは
    「〜である」と言い切る。断定できないなら、なぜ断定できないかを書く。
  \item[主観語を使わない] 「非常に」「かなり」「劇的に」は情報を持たない。
    「3 倍大きい」と書く。
  \item[「思う」は使わない] 論文は感想文ではない。
  \item[用語を統一する] 「対流セル」と「対流構造」を無意識に混ぜない。
    同じものは同じ語で呼び、違うものは違う語で呼ぶ。
  \item[数式も文の一部] 式の前後の助詞をつなげて音読し、
    日本語として成立するか確かめる。
\end{description}

\subsection{単位と有効数字}

物理量には必ず単位をつける。数値の有効数字は測定・計算の精度に見合った桁数にする。
計算機が出した 8 桁をそのまま書き写さない。
単位つきの量は \texttt{siunitx} を使うと組版が揃う。

\begin{verbatim}
太陽半径は \qty{6.96e8}{m}、自転角速度は \qty{2.6e-6}{\per\second} である。
速度は \qty{2}{\km\per\second}、範囲は \qtyrange{1}{5}{\km\per\second}。
\end{verbatim}

\noindent
従来どおり \verb|$2~\mathrm{km~s^{-1}}$| と手で書いてもよいが、
論文全体でどちらかに統一すること。

\subsection{書き始められないとき}

完成形を頭の中で作ってから書こうとすると、いつまでも書き始められない。
図を先に作り、それを並べ替えて話の順番を決め、
各図に対して説明の段落を書く、という順で進めるとよい。
文章の質は推敲で上げるもので、初稿で上げるものではない。


\section{\LaTeX{} の使い方}
\label{sec:guide_latex}

\subsection{このテンプレートの構成}

\begin{description}
  \item[\texttt{main.tex}] プリアンブル、表紙に載る情報、概要、章の並び。
    自分用に書き換えるのは主にこのファイルの上半分である。
  \item[\texttt{body.tex}] 本文。序論から結論まで。
  \item[\texttt{guide.tex}] この手引き。書き終わったら外してよい。
  \item[\texttt{thesis.bib}] 引用文献のデータベース。
  \item[\texttt{achilles.sty}] 全テンプレート共通のプリアンブル。
    原本はリポジトリの \texttt{common/} にある。
  \item[\texttt{thesiscover.sty}] 表紙のレイアウト。通常は触らない。
  \item[\texttt{fig/}] 図を置くところ。
\end{description}

\subsection{コンパイル}

このディレクトリで
\begin{verbatim}
latexmk
\end{verbatim}
とすれば \texttt{main.pdf} ができる。書きながら自動で作り直したいときは
\begin{verbatim}
latexmk -pvc
\end{verbatim}
とする。中間ファイルがおかしくなったら \verb|latexmk -C| で全部消してやり直す。
エンジンは LuaLaTeX である。\texttt{platex} や \texttt{uplatex} では通らない。

\subsection{相互参照}

図・表・式・章には \verb|\label{}| をつけ、本文からは \verb|\ref{}| で参照する。
番号を手で書いてはいけない。あとで章を入れ替えたときに必ず破綻する。
ラベルには \verb|fig:|、\verb|tab:|、\verb|eq:|、\verb|sec:|、\verb|chap:| の
接頭辞をつけると、何を指しているか分かりやすい。

\begin{verbatim}
図\ref{fig:sample}に示すように……
式\eqref{eq:induction}より……
\end{verbatim}

\noindent
式の参照は \verb|\ref| ではなく \verb|\eqref| を使うと括弧がつく。

\subsection{引用}

\texttt{thesis.bib} に文献情報を書き、本文から \verb|\citep| か \verb|\citet| で引く。
NASA ADS の \textsf{Export Citation} から \textsf{BibTeX} を選べば、
そのまま貼り付けられる形式で得られる。手で打ち直さないこと。

\begin{verbatim}
小スケールダイナモが有効粘性を下げる \citep{hotta2021}。
\citet{hotta2021} はこれを高解像度計算で示した。
\end{verbatim}

\noindent
このテンプレートは \texttt{biblatex + biber} と \texttt{natbib + bst} の
どちらでも動くようにしてある。\texttt{main.tex} の冒頭で切り替える。
本文の書き方は両者で共通なので、あとから変えても原稿は書き直さなくてよい。

\subsection{数式}

番号をつけて後で参照する式は \texttt{align}、参照しない式は \texttt{align*} を使う。
\verb|$$ ... $$| は使わない（\LaTeX{} では非推奨で、前後の空きがおかしくなる）。
ベクトルは \verb|\vect{B}| と書くと太字斜体になる。
\verb|\grad|、\verb|\divergence|、\verb|\curl|、\verb|\pdif{}{}| などの
短縮マクロを \texttt{achilles.sty} に用意してある。

\subsection{よくあるつまずき}

\begin{description}
  \item[エラーが出たら一番上から読む] 後ろのエラーは前のエラーの巻き添えである
    ことが多い。\verb|! LaTeX Error:| で始まる最初の行と、その直前の行番号を見る。
  \item[Undefined control sequence] マクロの綴りが違うか、
    必要なパッケージを読み込んでいない。
  \item[Overfull \textbackslash hbox] 行がはみ出している。長い英単語や URL が原因。
    数 pt なら無視してよいが、目で見て飛び出していたら直す。
  \item[番号が \texttt{??} になる] 参照が解決していない。もう一度コンパイルする
    （\texttt{latexmk} なら自動でやる）。それでも直らないならラベルの綴り間違い。
  \item[図が出ない] ファイル名か拡張子が違う。LuaLaTeX が読めるのは
    PDF・PNG・JPEG である。EPS は事前に PDF へ変換しておく。
\end{description}


\section{図表の作り方}
\label{sec:guide_figure}

審査員は本文より先に図を見る。図が分かりやすければ論文は半分伝わったも同然である。
逆に、図が読めない論文は本文がどれだけ良くても評価されない。

\subsection{形式}

\begin{description}
  \item[線画はベクタ形式で] 折れ線・等高線・散布図は PDF で保存する。
    拡大しても線が滑らかで、印刷しても綺麗である。
    Matplotlib なら \verb|plt.savefig('foo.pdf')| とするだけである。
  \item[画像はラスタ形式で] 太陽表面の強度分布のような画素の塊は PNG にする。
    PDF にすると数十 MB になり、PDF ビューアが固まる。
    \verb|dpi=300| 程度あれば印刷に耐える。
  \item[JPEG は使わない] 圧縮による滲みが出る。写真以外では使わない。
\end{description}

\subsection{文字の大きさ}

図の中の文字が本文より小さいと読めない。
論文に貼ったときに\textbf{本文と同じかやや小さい程度}になるよう、
図を作る段階でフォントサイズを決める。
\verb|\includegraphics[width=0.5\linewidth]| のように後から縮めると
文字も一緒に縮むので、その分を見越しておく。
出来上がった PDF を 100\% 表示にして、読めるか自分の目で確かめること。

\subsection{色}

\begin{description}
  \item[虹色（jet）を使わない] 明度が単調でないため、
    データにない偽の構造が見えてしまう。Matplotlib の既定である
    \texttt{viridis} や、正負のある量には \texttt{RdBu\_r} のような
    発散型カラーマップを使う。
  \item[色覚多様性に配慮する] 日本人男性の約 5\% は赤と緑の区別がつきにくい。
    赤と緑だけで系列を区別しない。線種（実線・破線・点線）や
    マーカーの形も併せて変える。
  \item[白黒で印刷されても読めるか] 製本された論文が白黒コピーされることはある。
    色を消しても情報が失われないようにしておくと安全である。
  \item[色に意味を持たせたら統一する] 論文全体で「観測は黒、計算は赤」と決めたら、
    最後まで変えない。
\end{description}

\subsection{キャプション}

キャプションだけを読んで図が理解できるように書く。含めるべきものは
\begin{enumerate}
  \item 何を示した図か
  \item 軸の量と単位
  \item 線・色・記号がそれぞれ何を表すか
  \item どの計算・観測のデータか
\end{enumerate}
である。「結果」だけのキャプションは論外である。
図のキャプションは図の\textbf{下}、表のキャプションは表の\textbf{上}に置く。

\subsection{図を並べる}

\texttt{subcaption} パッケージで (a) (b) と番号を振れる。

\begin{verbatim}
\begin{figure}[tb]
  \centering
  \begin{subfigure}{0.48\linewidth}
    \includegraphics[width=\linewidth]{left.pdf}
    \caption{低解像度}\label{fig:comp_low}
  \end{subfigure}\hfill
  \begin{subfigure}{0.48\linewidth}
    \includegraphics[width=\linewidth]{right.pdf}
    \caption{高解像度}\label{fig:comp_high}
  \end{subfigure}
  \caption{解像度による対流構造の違い。}
  \label{fig:comp}
\end{figure}
\end{verbatim}

\subsection{図の位置が思いどおりにならないとき}

\LaTeX{} は図を「浮きもの」として扱い、ページの切れ目を見て適当な場所に置く。
\verb|[tb]| は「ページの上か下に置いてよい」という指定である。
\verb|[h]|（ここに置け）を多用すると、かえって図が後ろに溜まる。
最終稿になるまでは位置を気にせず、内容を書くことに集中したほうがよい。


\section{引用と剽窃}
\label{sec:guide_ethics}

これは作法の問題ではなく、研究者としての資格の問題である。
剽窃が発覚すれば学位は取り消される。指導教員にも累が及ぶ。
知らなかったでは済まないので、必ず読むこと。

\subsection{何が剽窃か}

他人の文章・図・アイデアを、出典を示さずに自分のものとして提示することは
すべて剽窃である。以下はいずれも剽窃にあたる。

\begin{itemize}
  \item 論文や web の文章をコピーして貼り、少し語尾を変えて自分の文にする
  \item 他人の図をキャプションに出典を書かずに載せる
  \item 先輩の修士論文から章をまるごと持ってくる
  \item 生成 AI が出力した文章を、出典も検証もなくそのまま載せる
  \item 自分の過去の論文から、断りなく文章を再利用する（自己剽窃）
\end{itemize}

\subsection{正しいやり方}

\begin{description}
  \item[事実や知見を借りるとき] 自分の言葉で書き直し、引用をつける。
    「対流層の底は $0.71\,R_\odot$ にある \citep{christensen-dalsgaard1991}。」
  \item[文章をそのまま使うとき] 引用符でくくり、出典を示す。
    ただし自然科学の論文で原文引用が必要になることはまずない。
    必要だと思ったら、たいていは自分の言葉で書き直すべき場面である。
  \item[図を借りるとき] キャプションの末尾に出典を書く。
    「\citet{hotta2021} の Figure 3 を改変。」
    商業出版社の論文図は、多くの場合、著者本人でも転載に許諾が要る。
    論文の web ページにある \textsf{Permissions} や \textsf{RightsLink} から
    手続きする。許諾が取れないなら、自分で描き直す。
\end{description}

\noindent
どこまで自分の言葉にすればよいかの基準は、
語順を入れ替えたか同義語に置き換えたかではない。
\textbf{元の文章を見ずに、理解した内容を自分で書けるか}である。
見ないと書けないなら、まだ理解していない。

\subsection{生成 AI の扱い}

英文の校正や、書き出しに詰まったときの相談に使うのは構わない。
ただし次の 3 点は守ること。

\begin{enumerate}
  \item 出力された内容の正しさは自分で確認する。
    存在しない論文を引用してくることがある。
  \item 引用文献は必ず ADS などの一次情報で実在を確かめる。
  \item 使い方について、指導教員に隠さない。
    大学や学会が定める方針にも従うこと。
\end{enumerate}

\noindent
自分が説明できない文章を論文に載せてはいけない。
審査で問われるのは論文ではなく、書いた本人だからである。

\subsection{データと図の加工}

都合の悪い点を消す、エラーバーを小さく見せる、
軸を切って差を大きく見せる、といった操作は捏造・改竄である。
処理を加えたなら、加えたことを本文に書く。
「見やすさのために平滑化した」と書けば済む話を、黙ってやると不正になる。


\section{原稿の管理と添削の受け方}
\label{sec:guide_review}

\subsection{Git で履歴を残す}

\texttt{.tex} はテキストファイルなので Git で管理できる。
締切前に「昨日の版に戻したい」と言い出す前に、最初から履歴を残しておくこと。
節を書き終えるたび、あるいは一日の終わりに
\begin{verbatim}
git add -A
git commit -m "序論の先行研究を追加"
\end{verbatim}
としておけばよい。中間ファイルや PDF は \texttt{.gitignore} で除外してある。
バックアップは Git だけに頼らず、クラウド（大学のストレージなど）にも置くこと。

\subsection{添削に出す前に自分でやること}

指導教員の時間を、自分で見つけられる誤りに使わせないこと。
出す前に次を済ませておく。

\begin{itemize}
  \item \verb|\todo| と \verb|\red| が残っていないか検索する
  \item \texttt{??} や \texttt{[?]} が PDF に出ていないか確認する
  \item 図がすべて本文から参照されているか確認する
  \item 通しで一度、声に出して読む（黙読よりずっと効く。文の切れ目が分かる）
  \item 数式の記号がすべて本文で定義されているか確認する
\end{itemize}

\subsection{添削の受け方}

\begin{description}
  \item[早く出す] 完璧になってから見せようとすると間に合わない。
    構成だけの段階、章ごと、と何度も出したほうが結果的に速い。
  \item[版に名前をつける] PDF は \texttt{thesis\_20270105.pdf} のように
    日付をつけて渡す。「最新版」が何を指すか分からなくなるのを防ぐ。
  \item[指摘の理由を理解する] 直せと言われた箇所だけ直すと、
    同じ誤りを他の場所に残したまま次の版を出すことになる。
    なぜそう指摘されたのかを考え、同種の箇所をまとめて直す。
  \item[納得できないときは聞く] 指摘が間違っていることもある。
    黙って従うのも、黙って無視するのも良くない。
  \item[直したかどうかを返す] どの指摘をどう直したか、
    直さなかったならなぜかを添えて返すと、次の添削が速くなる。
\end{description}

\subsection{変更点を見せる}

\texttt{latexdiff} を使うと、前の版からの変更箇所に下線と取り消し線がついた
PDF を作れる。どこを直したかが一目で分かるので、
再添削のときはこれを添えると親切である。

\begin{verbatim}
latexdiff-vc --git -r HEAD~1 --flatten main.tex
\end{verbatim}

\subsection{提出前の最終確認}

\begin{itemize}
  \item 表紙の年度・題目・氏名・学籍番号・所属が学務の書類と一致しているか
  \item ページ番号が通っているか
  \item 目次・図目次・表目次の番号が本文と合っているか
  \item 引用文献が本文の引用とすべて対応しているか
  \item 手引き（この章）を \texttt{main.tex} から外したか
  \item 行番号（\texttt{lineno}）を切ったか
  \item 印刷して提出するなら、リンクの色を黒にしたか
  \item PDF を実際に印刷して、図の文字が読めるか確かめたか
\end{itemize}


%% 引用文献。main.tex の末尾と同じものです。
%% どちらの方式を使うかは preamble.tex の冒頭で決めています。
\ifdefined\BibStyleBiblatex
  \printbibliography[title={引用文献}]
\else
  \renewcommand{\bibname}{引用文献}
  \bibliographystyle{plainnat}
  \bibliography{thesis}
\fi

\end{document}
